\section*{ID8108: Example (Dipole Near Field) — Quick Symbolic Sketch (1)}
\paragraph{Metadata.}
PiID: 4070866711496 | RTEID: RTE:P31485.1646:T1.6124 | UUID: 184fbdea-8132-5487-911b-f964605cf8c1

\paragraph{Canonical Statement.}
\# 6) Example (dipole near field) — quick symbolic sketch Given your dipole vector potential (in the file) \textbackslash{}[ A\_\{\textbackslash{}rm dipole\}(\textbackslash{}mathbf r)=\textbackslash{}frac\{\textbackslash{}mu\_0\}\{4\textbackslash{}pi\}\textbackslash{}frac\{\textbackslash{}boldsymbol\{\textbackslash{}mu\}\_\{\textbackslash{}rm eff\}\textbackslash{}times(\textbackslash{}mathbf r-\textbackslash{}mathbf r\_\{\textbackslash{}rm cm\})\}\{|\textbackslash{}mathbf r-\textbackslash{}mathbf r\_\{\textbackslash{}rm cm\}|\textasciicircum{}3\}, \textbackslash{}] compute \textbackslash{}[ B=\textbackslash{}nabla\textbackslash{}times A\_\{\textbackslash{}rm dipole\},\textbackslash{}qquad E=-\textbackslash{}partial\_t A - \textbackslash{}nabla\textbackslash{}Phi. \textbackslash{}] Then build \textbackslash{}[ T\textasciicircum{}\{00\}\_\{\textbackslash{}rm EM\} \textbackslash{}simeq \textbackslash{}mathcal\{E\}\_\{\textbackslash{}rm EM\},\textbackslash{}qquad T\textasciicircum{}\{0i\}\_\{\textbackslash{}rm EM\}\textbackslash{}simeq \textbackslash{}frac\{1\}\{c\}S\textasciicircum{}i\_\{\textbackslash{}rm EM\},\textbackslash{}qquad T\textasciicircum{}\{ij\}\_\{\textbackslash{}rm EM\}=\textbackslash{}sigma\textasciicircum{}\{ij\}. \textbackslash{}] If the dipole is slowly time-varying and the rotation is nonrelativistic (\textbackslash{}(|\textbackslash{}beta|\textbackslash{}ll1\textbackslash{})), you can use these 3-vector forms directly and then correct order-by-order in \textbackslash{}(\textbackslash{}beta=(\textbackslash{}omega\textbackslash{}times r)/c\textbackslash{}) per your Lorentz transform for the rotating frame. Your file already includes the Lorentz transformation pieces (the \textbackslash{}(\textbackslash{}Lambda\textbackslash{}) matrix and \textbackslash{}(\textbackslash{}beta\textbackslash{}), \textbackslash{}(\textbackslash{}gamma\textbackslash{})). fileciteturn0file0

\paragraph{Canonical Equation.}
\[
B=\nabla\times A_{\rm dipole},\qquad E=-\partial_t A - \nabla\Phi.
\]

\paragraph{Dictionary.}
Relation: B=∇× A\_rm dipole, E=-∂\_t A - ∇Φ.

\paragraph{Encyclopedia.}
Example (Dipole Near Field) — Quick Symbolic Sketch (1) is a differential equation, written B=∇× A\_rm dipole, E=-∂\_t A - ∇Φ.. It relates the quantity to its derivatives, governing how the system evolves in space and time. It is built from ∇, Φ, A\_ dipole, E, A, defining B. Within the Trawin Zero Point registry it serves as one element of the framework's mathematical narrative.
